% GENERATED FILE. Edit canonical lessons, then run npm run generate:books.
% canonical-source-hash: fnv1a64:f8c2be74c8c1688e
% canonical-lessons: MR-C270-badam, MR-C270-manuka, MR-C270-gul, MR-C270-dhotar, MR-C270-kurta

\chapter{Almonds, Raisins, Jaggery, Dhoti, Kurta}
\label{ch:mr-almonds-raisins-jaggery-dhoti-kurta}
\hlchaptermodality{270}

\begin{chapteropening}
\textbf{By the end of this chapter:} \emph{I can say almonds, raisins, jaggery, a dhoti and a kurta.}

Complete the last lesson of chapter 270: I can say almonds, raisins, jaggery, a dhoti and a kurta.
\end{chapteropening}

\section[\texorpdfstring{\mr{बदाम} (badām)}{बदाम (badām)}]{\mr{बदाम} (badām) — almonds}

\label{lesson:MR-C270-badam}



\begin{quote}
Before the new one: say the Marathi for a sprouted-bean curry, then the Marathi for a pickle.
\end{quote}

\subsection*{You'll want to know: \mr{बदाम}}
\textbf{\mr{बदाम}} — \emph{badām} — \textquotedblleft{}almonds\textquotedblright{}.

\begin{grammarlens}[title={What you've built}]
One more word for this chapter.
\end{grammarlens}

\subsection*{Your turn}
\begin{itemize}
\raggedright
  \item \emph{Say it:} \emph{badām}
  \item \emph{Say it:} \emph{badām}, once more
  \item \emph{Say it:} say \emph{badām}
  \item \emph{Recall:} say the Marathi for a sprouted-bean curry, then the Marathi for a pickle, then say \emph{badām} again
\end{itemize}

\begin{tcolorbox}[breakable,colback=teal!4,colframe=teal!35!black,title={Before you move on}]
What does \mr{बदाम} mean? (\textquotedblleft{}Almonds\textquotedblright{}.) Say it once more.
\end{tcolorbox}
\section[\texorpdfstring{\mr{मनुका} (manukā)}{मनुका (manukā)}]{\mr{मनुका} (manukā) — raisins}

\label{lesson:MR-C270-manuka}



\begin{quote}
Before the new one: say the Marathi for a pickle, then the Marathi for almonds.
\end{quote}

\subsection*{You'll want to know: \mr{मनुका}}
\textbf{\mr{मनुका}} — \emph{manukā} — \textquotedblleft{}raisins\textquotedblright{}.

\begin{grammarlens}[title={What you've built}]
One more word for this chapter.
\end{grammarlens}

\subsection*{Your turn}
\begin{itemize}
\raggedright
  \item \emph{Say it:} \emph{manukā}
  \item \emph{Say it:} \emph{manukā}, once more
  \item \emph{Say it:} say \emph{manukā}
  \item \emph{Recall:} say the Marathi for a pickle, then the Marathi for almonds, then say \emph{manukā} again
\end{itemize}

\begin{tcolorbox}[breakable,colback=teal!4,colframe=teal!35!black,title={Before you move on}]
What does \mr{मनुका} mean? (\textquotedblleft{}Raisins\textquotedblright{}.) Say it once more.
\end{tcolorbox}
\section[\texorpdfstring{\mr{गूळ} (gūḷ)}{गूळ (gūḷ)}]{\mr{गूळ} (gūḷ) — jaggery}

\label{lesson:MR-C270-gul}



\begin{quote}
Before the new one: say the Marathi for almonds, then the Marathi for raisins.
\end{quote}

\subsection*{You'll want to know: \mr{गूळ}}
\textbf{\mr{गूळ}} — \emph{gūḷ} — \textquotedblleft{}jaggery\textquotedblright{}.

\begin{grammarlens}[title={What you've built}]
One more word for this chapter.
\end{grammarlens}

\subsection*{Your turn}
\begin{itemize}
\raggedright
  \item \emph{Say it:} \emph{gūḷ}
  \item \emph{Say it:} \emph{gūḷ}, once more
  \item \emph{Say it:} say \emph{gūḷ}
  \item \emph{Recall:} say the Marathi for almonds, then the Marathi for raisins, then say \emph{gūḷ} again
\end{itemize}

\begin{tcolorbox}[breakable,colback=teal!4,colframe=teal!35!black,title={Before you move on}]
What does \mr{गूळ} mean? (\textquotedblleft{}Jaggery\textquotedblright{}.) Say it once more.
\end{tcolorbox}
\section[\texorpdfstring{\mr{धोतर} (dhotar)}{धोतर (dhotar)}]{\mr{धोतर} (dhotar) — a dhoti}

\label{lesson:MR-C270-dhotar}



\begin{quote}
Before the new one: say the Marathi for raisins, then the Marathi for jaggery.
\end{quote}

\subsection*{You'll want to know: \mr{धोतर}}
\textbf{\mr{धोतर}} — \emph{dhotar} — \textquotedblleft{}a dhoti\textquotedblright{}.

\begin{grammarlens}[title={What you've built}]
One more word for this chapter.
\end{grammarlens}

\subsection*{Your turn}
\begin{itemize}
\raggedright
  \item \emph{Say it:} \emph{dhotar}
  \item \emph{Say it:} \emph{dhotar}, once more
  \item \emph{Say it:} say \emph{dhotar}
  \item \emph{Recall:} say the Marathi for raisins, then the Marathi for jaggery, then say \emph{dhotar} again
\end{itemize}

\begin{tcolorbox}[breakable,colback=teal!4,colframe=teal!35!black,title={Before you move on}]
What does \mr{धोतर} mean? (\textquotedblleft{}A dhoti\textquotedblright{}.) Say it once more.
\end{tcolorbox}
\section[\texorpdfstring{\mr{कुर्ता} (kurtā)}{कुर्ता (kurtā)}]{\mr{कुर्ता} (kurtā) — a kurta}

\label{lesson:MR-C270-kurta}



\begin{quote}
Before the new one: say the Marathi for jaggery, then the Marathi for a dhoti.
\end{quote}

\subsection*{You'll want to know: \mr{कुर्ता}}
\textbf{\mr{कुर्ता}} — \emph{kurtā} — \textquotedblleft{}a kurta\textquotedblright{}.

\begin{grammarlens}[title={What you've built}]
One more word for this chapter.
\end{grammarlens}

\subsection*{Your turn}
\begin{itemize}
\raggedright
  \item \emph{Say it:} \emph{kurtā}
  \item \emph{Say it:} \emph{kurtā}, once more
  \item \emph{Say it:} say \emph{kurtā}
  \item \emph{Recall:} say the Marathi for jaggery, then the Marathi for a dhoti, then say \emph{kurtā} again
\end{itemize}

\begin{tcolorbox}[breakable,colback=teal!4,colframe=teal!35!black,title={Before you move on}]
What does \mr{कुर्ता} mean? (\textquotedblleft{}A kurta\textquotedblright{}.) Say it once more.
\end{tcolorbox}
